\section{tabularray 集成与单元格标签对齐}
\pkg{tabularray} 适配库需要 2025C（2025-11-27）或更新版本。
在导言区加载，普通 \pkg{gradbars} 文档不必安装或加载此可选依赖：
\begin{codeonly}
\usepackage{tabularray}
\usepackage{gradbars}
\UseTblrLibrary{gradbars}
\end{codeonly}
库文件是 \key{tblrlibgradbars.sty}，与其余运行文件一起放在 \key{sty/}。
它按需加载 Heat 和 \pkg{tabularray} 的 hook/varwidth 支持；重复加载无副作用。
只加载 \pkg{tabularray} 再调用上述库入口时，也会自动加载最小的 \pkg{gradbars} 核心。

\subsection{专用列与单元格命令}
专用列采用 \key{G[选项]{宽度}}、\key{H[选项]{宽度}}、
\key{N[选项]{宽度}} 和 \key{D[选项]{宽度}}，分别绘制条形、Heat、
数字和双层比较条。\key{D} 的数据为 \key{参考值/当前值}；整格为空或 \key{NA} 时表示两个值均缺失。
宽度指内容区域，不包含 \key{colsep} 或竖线；表头用一对花括号保护，避免被当作数值处理。
这些列名由本库注册，已有同名自定义列应先改名。
\begin{example}[text and listing,sidebyside=false]{gradtblr 预设与正负单元格}
\begingroup
\gradtblrsetup{min=-10,max=10,precision=1,height=14pt,
  color=gradbarsTeal!45,negative color=red!35,
  track color=white,text color=black}
\begin{gradtblr}{colspec={l G{23mm} H{23mm} N{23mm} D{35mm}}}
Method & {Change} & {Heat} & {Number} & {Before/after}\\
A & 0.5 & -5 & 1.25 & -3/7\\
B & -9.5 & 5 & -2.5 & 4/NA\\
Pending & NA & NA & NA & NA/NA\\
\end{gradtblr}
\endgroup
\end{example}
\key{gradtblr} 是提供居中单元格、浅色粗体表头、细横线及适当留白的独立环境；
\key{tblr} 本身的默认样式不会改变。表头仍须保护。
\key{\gradtblrsetup} 在当前 TeX 分组内累积适配器默认设置。
覆盖顺序为普通图形设置、适配器初始设置、\key{\gradtblrsetup}、当前调用选项、最终分配宽度。
适配器初始设置是 \key{baseline=text,label=inside,label align=cell}。

现有 \key{Q} 或 \key{X} 列可通过 \key{cmd} 选择对应命令：
\begin{codeonly}
cmd=\gradbarTblr       % 条形
cmd=\gradheatTblr      % Heat
cmd=\gradnumberTblr    % 数字
cmd=\gradcompareTblr   % 双层比较条
\end{codeonly}
直接调用时，在命令后添加可选设置和数据参数。
通用单元格命令的语法为：
\begin{codeonly}
\gradtblrcell[选项]{类型}{数据}
\end{codeonly}
支持的类型为
\key{bar}、\key{heat}、\key{number}、\key{compare}、\key{dumbbell} 和 \key{range}；
后三种类型均接受两个斜线分隔值，\key{range} 为下界/上界。
数字列复用精度、格式化器和单位，但不裁剪数值，也不创建图形替代文本。
\pkg{siunitx} 仍须由用户显式加载。

\subsection{以单元格为基准居中及左侧标签}
\key{label align=cell} 配合 \key{label=inside} 或 \key{label=auto}，
使标签位于完整图形宽度的中心，不跟随当前填充段、正负方向或标签长度改变位置。
\key{label align=graphic} 保留原有图形对齐规则，且是普通图形的默认值。
\key{label=none}、\key{label=left}、\key{label=outside} 和 \key{label=end} 保留各自位置语义。
\key{label overflow=allow} 仅允许越界，不选择居中基准。

\key{cell width=auto} 使用当前内容宽度；自然列初次测量以原图形宽度为估计，
随后使用已分配的宽度。固定列、\key{X} 列和跨列单元格均可使用；
在表格外直接调用时，建议显式指定 \key{cell width}。
显式宽度不得超过表格实际分配的内容区域。
左侧/外侧标签会从总宽度中预留 \key{label width} 与 \key{label gap}。
文字过长时应增加槽宽、内容宽度或减小字体；默认报错，\key{label overflow=allow} 可允许有意越界。
\begin{example}[text and listing,sidebyside=false]{左侧百分号与单元格宽度}
\begingroup
\gradtblrsetup{min=0,max=.065,precision=5,unit=\%,
  label=left,label width=16mm,height=8pt}
\begin{tblr}{colspec={l G{45mm}},row{1}={font=\bfseries}}
Decade & {Frequency}\\
1700 & .05942\\
1710 & .03044\\
\end{tblr}
\endgroup
\end{example}

\subsection{跨页与无障碍边界}
\key{longtblr} 可复用上述列，用 \key{rowhead=1} 重复受保护的文本表头。
请直接放在单栏正文中，避免不可分页容器：
\begin{codeonly}
\begin{longtblr}{colspec={l G[min=0,max=100]{40mm}},rowhead=1}
Method & {Score}\\
A & 72\\
B & 94\\
\end{longtblr}
\end{codeonly}
本适配层负责单元格图形，不接管 CSV 排序、自动范围、排名和汇总；
需要这些流程时采用 \key{gradtable}，需要自由合并与 \pkg{tabularray} 布局时采用本库。
显式 \key{alt}、自动描述及 \key{accessibility=artifact|off} 沿用图形语义。
测量阶段暂停标记，避免重复替代文本；重复表头应使用文字，含图形表头的逐页替代文本不在当前验证范围。
图形描述支持不等于 \pkg{tabularray} 表格结构或整份文档满足 PDF/UA。
